\begin{mvtable}
\centering
{\zihao{6}续表\par}\vspace{2pt}
\begin{tabular}{cccc}
\toprule
\multirow{2}{*}{样本容量 $n$} & \multicolumn{3}{c}{显著性水平 $\alpha$} \\
\cmidrule(lr){2-4}
 & 0.01 & 0.05 & 0.10 \\
\midrule
50 & 0.967 1 & 0.976 8 & 0.980 9 \\
55 & 0.969 5 & 0.978 7 & 0.982 2 \\
60 & 0.972 0 & 0.980 1 & 0.983 6 \\
75 & 0.977 1 & 0.983 8 & 0.986 6 \\
100 & 0.982 2 & 0.987 3 & 0.989 5 \\
150 & 0.987 9 & 0.991 3 & 0.992 8 \\
200 & 0.990 5 & 0.993 1 & 0.994 2 \\
300 & 0.993 5 & 0.995 3 & 0.996 0 \\
\bottomrule
\end{tabular}
\end{mvtable}

\noindent\textbf{例 4.11}\quad \textbf{(正态性的相关系数检验法)}

让我们计算例 4.9(参见\hyperref[fig:4.5]{图 4.5})的 Q-Q 图的相关系数 $r_{Q}$ 并检验正态性.

利用例 4.9 的数据, 我们有 $\bar{x}=0.770$ 以及
\[
\sum_{j=1}^{10}(x_{(j)}-\bar{x})q_{(j)}=8.584,\quad \sum_{j=1}^{10}(x_{(j)}-\bar{x})^{2}=8.472\quad\text{和}\quad \sum_{j=1}^{10}q_{(j)}^{2}=8.795
\]

由于 $\bar{q}=0$
\[
r_{Q}=\frac{8.584}{\sqrt{8.472\times 8.795}}=0.994
\]

在显著性水平 10\% 下, 通过把 $r_{Q}=0.994$ 与\hyperref[tab:4.2]{表 4.2} 中对应于 $n=10$ 和 $\alpha=0.10$ 的表项相对照, 来进行正态性检验. 表项为 0.935 1. 由于 $r_{Q}>0.935\ 1$, 我们不拒绝正态性假定.\hfill$\blacksquare$

与 $r_{Q}$ 不同, 有些软件包评估由夏皮罗(Shapiro)和威尔克(Wilk)(参见文献[12])提出的原始统计量. 它的相关形式对应于用一个标准正态顺序统计量的期望值和它们的协方差的函数来替换 $q_{(j)}$. 我们喜欢 $r_{Q}$, 因为它直接对应于正态的许多散布图中的点. 对于大样本容量, 这两个统计量几乎是相同的(参见文献[13]), 因此, 任何一个都可用来判断拟合的不足.

可以研究多于一个特性的线性组合. 许多统计学家建议标绘
\[
\hat{\mathbf{e}}_{1}'\mathbf{x},\quad \text{其中}\quad \mathbf{S}\hat{\mathbf{e}}_{1}=\hat{\lambda}_{1}\hat{\mathbf{e}}_{1}
\]
式中的 $\hat{\lambda}_{1}$ 是 $\mathbf{S}$ 的最大特征值. 这里 $\mathbf{x}_{(j)}'=[x_{(j1)},x_{(j2)},\cdots,x_{(jp)}]$ 是对 $p$ 个变量 $X_{1},X_{2},\cdots,X_{p}$ 的第 $j$ 个的观测值. 对应于最小特征值的线性组合 $\hat{\mathbf{e}}_{p}'\mathbf{x}_{(j)}$ 也常常被挑选出来进行检验(有关进一步的细节参见第 8 章和文献[6]).

\mvsub{评估二元正态性}

我们会希望检验所有 $2,3,\cdots,p$ 维的分布的正态性假定. 然而, 正如我们已指出的, 对于实际工作, 研究一元和二元分布通常就足够了. 前面我们考虑过一元边缘分布. 现在感兴趣的是检查二元情况.

在第 1 章中, 我们对于表示特征的点对作了散布图. 如果观测值是从一个多元正态分布生成的, 每个二元分布会是正态的, 并且常数密度轮廓线会是椭圆. 散布图显示一个近乎椭圆
